% GENERATED FILE. Edit canonical lessons, then run npm run generate:books.
% canonical-source-hash: fnv1a64:321ed3811527b338
% canonical-lessons: TA-C477-marataippu, TA-C477-purrunoy, TA-C477-pakkavatam, TA-C477-astuma, TA-C477-kamalai

\chapter{Serious Illness}
\label{ch:ta-serious-illness}
\hlchaptermodality{477}

\begin{chapteropening}
\textbf{By the end of this chapter:} \emph{I can name serious illnesses such as heart attacks, cancer and strokes.}

Complete the last lesson of chapter 477: I can name serious illnesses such as heart attacks, cancer and strokes.
\end{chapteropening}

\section[\texorpdfstring{\ta{மாரடைப்பு} (māraṭaippu)}{மாரடைப்பு (māraṭaippu)}]{\ta{மாரடைப்பு} (māraṭaippu) — a heart attack}

\label{lesson:TA-C477-marataippu}



\begin{quote}
Before the new one: say the Tamil for hygiene, then the Tamil for to fall ill.
\end{quote}

\subsection*{You'll want to know: \ta{மாரடைப்பு}}
\textbf{\ta{மாரடைப்பு}} — \emph{māraṭaippu} — \textquotedblleft{}a heart attack\textquotedblright{}.

\begin{quote}
\ta{தாத்தாவுக்கு} \ta{மாரடைப்பு} \ta{வந்ததால்} \ta{மருத்துவமனைக்குப்} \ta{போனோம்}.
\end{quote}

\begin{quote}
\emph{tāttāvukku māraṭaippu vantatāl maruttuvamaṉaikkup pōṉōm.}
\end{quote}

\begin{quote}
We went to the hospital because grandfather had a heart attack.
\end{quote}

\ta{மார்} \textquotedblleft{}chest\textquotedblright{} + \emph{aṭaippu} \textquotedblleft{}a blockage\textquotedblright{}: the heart's vessels close.

\begin{cousinweb}
\textbf{\ta{மாரடைப்பு}} (\emph{māraṭaippu}) is a native compound built from \ta{மார்} \emph{mār}, \textquotedblleft{}breast, chest\textquotedblright{} (DEDR 4818), and \emph{aṭaippu}, \textquotedblleft{}a blockage\textquotedblright{}.
\end{cousinweb}

\begin{grammarlens}[title={What you've built}]
One more word for describing what you lived and felt.
\end{grammarlens}

\subsection*{Your turn}
\begin{itemize}
\raggedright
  \item \emph{Say it:} \emph{māraṭaippu}
  \item \emph{Say it:} \emph{māraṭaippu}, once more
  \item \emph{Say it:} \emph{tāttāvukku māraṭaippu vantatāl maruttuvamaṉaikkup pōṉōm.} — We went to the hospital because grandfather had a heart attack.
  \item \emph{Recall:} say the Tamil for hygiene, then the Tamil for to fall ill, then say \emph{māraṭaippu} again
\end{itemize}

\begin{tcolorbox}[breakable,colback=teal!4,colframe=teal!35!black,title={Before you move on}]
What does \ta{மாரடைப்பு} mean? (\textquotedblleft{}A heart attack\textquotedblright{}.) Say it once more.
\end{tcolorbox}
\section[\texorpdfstring{\ta{புற்றுநோய்} (puṟṟunōy)}{புற்றுநோய் (puṟṟunōy)}]{\ta{புற்றுநோய்} (puṟṟunōy) — cancer}

\label{lesson:TA-C477-purrunoy}



\begin{quote}
Before the new one: say the Tamil for to fall ill, then the Tamil for a heart attack.
\end{quote}

\subsection*{You'll want to know: \ta{புற்றுநோய்}}
\textbf{\ta{புற்றுநோய்}} — \emph{puṟṟunōy} — \textquotedblleft{}cancer\textquotedblright{}.

\begin{quote}
\ta{புற்றுநோய்க்கு} \ta{இப்போது} \ta{நல்ல} \ta{சிகிச்சை} \ta{இருக்கிறது}.
\end{quote}

\begin{quote}
\emph{puṟṟunōykku ippōtu nalla cikiccai irukkiṟatu.}
\end{quote}

\begin{quote}
There is good treatment for cancer now.
\end{quote}

\ta{புற்று} \textquotedblleft{}a cancerous growth\textquotedblright{} + \ta{நோய்} \textquotedblleft{}disease\textquotedblright{}.

\begin{cousinweb}
\textbf{\ta{புற்றுநோய்}} (\emph{puṟṟunōy}) is a native compound of \ta{புற்று} \emph{puṟṟu}, \textquotedblleft{}anything scurvy or cancerous\textquotedblright{} (DEDR 4336), and \ta{நோய்}, \textquotedblleft{}disease\textquotedblright{} (DEDR 3793).
\end{cousinweb}

\begin{grammarlens}[title={What you've built}]
One more word for describing what you lived and felt.
\end{grammarlens}

\subsection*{Your turn}
\begin{itemize}
\raggedright
  \item \emph{Say it:} \emph{puṟṟunōy}
  \item \emph{Say it:} \emph{puṟṟunōy}, once more
  \item \emph{Say it:} \emph{puṟṟunōykku ippōtu nalla cikiccai irukkiṟatu.} — There is good treatment for cancer now.
  \item \emph{Recall:} say the Tamil for to fall ill, then the Tamil for a heart attack, then say \emph{puṟṟunōy} again
\end{itemize}

\begin{tcolorbox}[breakable,colback=teal!4,colframe=teal!35!black,title={Before you move on}]
What does \ta{புற்றுநோய்} mean? (\textquotedblleft{}Cancer\textquotedblright{}.) Say it once more.
\end{tcolorbox}
\section[\texorpdfstring{\ta{பக்கவாதம்} (pakkavātam)}{பக்கவாதம் (pakkavātam)}]{\ta{பக்கவாதம்} (pakkavātam) — a stroke; paralysis}

\label{lesson:TA-C477-pakkavatam}



\begin{quote}
Before the new one: say the Tamil for a heart attack, then the Tamil for cancer.
\end{quote}

\subsection*{You'll want to know: \ta{பக்கவாதம்}}
\textbf{\ta{பக்கவாதம்}} — \emph{pakkavātam} — \textquotedblleft{}a stroke; paralysis\textquotedblright{}.

\begin{quote}
\ta{பக்கவாதம்} \ta{வந்த} \ta{பிறகு} \ta{அவரால்} \ta{நடக்க} \ta{முடியவில்லை}.
\end{quote}

\begin{quote}
\emph{pakkavātam vanta piṟaku avarāl naṭakka muṭiyavillai.}
\end{quote}

\begin{quote}
After the stroke he could not walk.
\end{quote}

\ta{பக்கம்} \textquotedblleft{}side\textquotedblright{} + \ta{வாதம்}, a disorder of wind in traditional medicine.

\begin{cousinweb}
\textbf{\ta{பக்கவாதம்}} (\emph{pakkavātam}) is borrowed from Sanskrit, built from \emph{pakṣa}, \textquotedblleft{}side\textquotedblright{} (CDIAL 7627), and \emph{vāta}, \textquotedblleft{}wind, a humour of the body\textquotedblright{} (CDIAL 11491).
\end{cousinweb}

\begin{grammarlens}[title={What you've built}]
One more word for describing what you lived and felt.
\end{grammarlens}

\subsection*{Your turn}
\begin{itemize}
\raggedright
  \item \emph{Say it:} \emph{pakkavātam}
  \item \emph{Say it:} \emph{pakkavātam}, once more
  \item \emph{Say it:} \emph{pakkavātam vanta piṟaku avarāl naṭakka muṭiyavillai.} — After the stroke he could not walk.
  \item \emph{Recall:} say the Tamil for a heart attack, then the Tamil for cancer, then say \emph{pakkavātam} again
\end{itemize}

\begin{tcolorbox}[breakable,colback=teal!4,colframe=teal!35!black,title={Before you move on}]
What does \ta{பக்கவாதம்} mean? (\textquotedblleft{}A stroke; paralysis\textquotedblright{}.) Say it once more.
\end{tcolorbox}
\section[\texorpdfstring{\ta{ஆஸ்துமா} (āstumā)}{ஆஸ்துமா (āstumā)}]{\ta{ஆஸ்துமா} (āstumā) — asthma}

\label{lesson:TA-C477-astuma}



\begin{quote}
Before the new one: say the Tamil for cancer, then the Tamil for a stroke.
\end{quote}

\subsection*{You'll want to know: \ta{ஆஸ்துமா}}
\textbf{\ta{ஆஸ்துமா}} — \emph{āstumā} — \textquotedblleft{}asthma\textquotedblright{}.

\begin{quote}
\ta{குளிர்} \ta{காலத்தில்} \ta{என்} \ta{தம்பிக்கு} \ta{ஆஸ்துமா} \ta{அதிகமாகும்}.
\end{quote}

\begin{quote}
\emph{kuḷir kālattil eṉ tampikku āstumā atikamākum.}
\end{quote}

\begin{quote}
In winter my younger brother’s asthma gets worse.
\end{quote}

Dative: \ta{அவனுக்கு} \ta{ஆஸ்துமா} \ta{இருக்கிறது}, he has asthma.

\begin{cousinweb}
\textbf{\ta{ஆஸ்துமா}} (\emph{āstumā}) is a loan from English \emph{asthma}, written with the grantha letter \ta{ஸ}.
\end{cousinweb}

\begin{grammarlens}[title={What you've built}]
One more word for describing what you lived and felt.
\end{grammarlens}

\subsection*{Your turn}
\begin{itemize}
\raggedright
  \item \emph{Say it:} \emph{āstumā}
  \item \emph{Say it:} \emph{āstumā}, once more
  \item \emph{Say it:} \emph{kuḷir kālattil eṉ tampikku āstumā atikamākum.} — In winter my younger brother’s asthma gets worse.
  \item \emph{Recall:} say the Tamil for cancer, then the Tamil for a stroke, then say \emph{āstumā} again
\end{itemize}

\begin{tcolorbox}[breakable,colback=teal!4,colframe=teal!35!black,title={Before you move on}]
What does \ta{ஆஸ்துமா} mean? (\textquotedblleft{}Asthma\textquotedblright{}.) Say it once more.
\end{tcolorbox}
\section[\texorpdfstring{\ta{காமாலை} (kāmālai)}{காமாலை (kāmālai)}]{\ta{காமாலை} (kāmālai) — jaundice}

\label{lesson:TA-C477-kamalai}



\begin{quote}
Before the new one: say the Tamil for a stroke, then the Tamil for asthma.
\end{quote}

\subsection*{You'll want to know: \ta{காமாலை}}
\textbf{\ta{காமாலை}} — \emph{kāmālai} — \textquotedblleft{}jaundice\textquotedblright{}.

\begin{quote}
\ta{காமாலை} \ta{வந்ததால்} \ta{அவனுடைய} \ta{கண்கள்} \ta{மஞ்சளாக} \ta{இருந்தன}.
\end{quote}

\begin{quote}
\emph{kāmālai vantatāl avaṉuṭaiya kaṇkaḷ mañcaḷāka iruntaṉa.}
\end{quote}

\begin{quote}
Because he had jaundice, his eyes were yellow.
\end{quote}

Often \ta{மஞ்சள்} \ta{காமாலை}, yellow jaundice.

\begin{cousinweb}
\textbf{\ta{காமாலை}} (\emph{kāmālai}) is borrowed from Sanskrit \emph{kāmalā}, \textquotedblleft{}jaundice\textquotedblright{} (CDIAL 3044).
\end{cousinweb}

\begin{grammarlens}[title={What you've built}]
One more word for describing what you lived and felt.
\end{grammarlens}

\subsection*{Your turn}
\begin{itemize}
\raggedright
  \item \emph{Say it:} \emph{kāmālai}
  \item \emph{Say it:} \emph{kāmālai}, once more
  \item \emph{Say it:} \emph{kāmālai vantatāl avaṉuṭaiya kaṇkaḷ mañcaḷāka iruntaṉa.} — Because he had jaundice, his eyes were yellow.
  \item \emph{Recall:} say the Tamil for a stroke, then the Tamil for asthma, then say \emph{kāmālai} again
\end{itemize}

\begin{tcolorbox}[breakable,colback=teal!4,colframe=teal!35!black,title={Before you move on}]
What does \ta{காமாலை} mean? (\textquotedblleft{}Jaundice\textquotedblright{}.) Say it once more.
\end{tcolorbox}
